\documentclass[11pt,a4paper]{article}
\usepackage[margin=1in]{geometry}
\usepackage{hyperref}
\usepackage{enumitem}
\usepackage{graphicx}
\usepackage{xcolor}
\usepackage{booktabs}

\makeatletter \makeatletter \newcommand{\BvrSetLabInstructor}[1]{%
  \gdef\Bvr@LabInstructor{#1}}

\newcommand{\Bvr@LabInstructor}{Dr. Raghav %
  B. Venkataramaiyer}

\newcommand{\BvrLabInstructorValue}{\Bvr@LabInstructor}
\makeatother

\usepackage{titling}

\pretitle{\begin{center}\bfseries\fontsize{18bp}{18bp}\selectfont}
\makeatletter
\newcommand{\BvrSubtitle}[1]{%
  \posttitle{%
    \par\end{center}
    \begin{center}\large#1\end{center}
    \vskip 0.5em}%
}
\makeatother

\newcommand{\BvrHint}[1]{%
  {\normalfont\normalsize\itshape\hfill #1}}


\title{Hostel Room Swapping Service}

\BvrSetLabInstructor{Dr. Raghav B. Venkataramaiyer}

\BvrSubtitle{%
  Project Proposal for UCS503P  \\
  Submitted to: \BvrLabInstructorValue \\ \bfseries
  Thapar Institute of Engineering and Technology%
}

\author{
  Aabhas Khandelwal, CSED%
  \thanks{Roll No: \texttt{1024030673}, %
    Email: \texttt{akhandelwal2\_be24\@thapar.edu}}
  \and
  Aakash Chawla, CSED%
  \thanks{Roll No: \texttt{1024031097}, %
    Email: \texttt{achawla1\_be24\@thapar.edu}}
  \and
  Joy Soni, CSED%
  \thanks{Roll No: \texttt{1024030674}, %
    Email: \texttt{jsoni\_be24\@thapar.edu}}
}

\date{\today}

\begin{document}
\maketitle

\section{Higher-order goal%
  \BvrHint{\ldots secondary framing}}
Make better use of hostel rooms that already exist.
Every semester students live in rooms that suit them
poorly while, in the same hostel, someone would gladly
take that room and give up theirs. Fixing this needs no
construction --- only visibility and coordination.

\section{Time-to-value %
\BvrHint{\ldots primary heuristic}}
\begin{itemize}[leftmargin=0.7cm]
    \item The matching engine was built and demonstrated
      end to end on synthetic data before any database
      existed, so the core proposition was validated in
      the first iteration.
    \item The domain layer has no framework or database
      dependency, so the full test suite runs in about
      two seconds on every push and a working
      demonstration is always available.
    \item Success is measured, not asserted: the gain in
      cohort preference satisfaction per matching round.
\end{itemize}

\section{Problem Statement}
Students are allotted rooms they did not choose --- wrong
floor, wrong direction, away from friends. They are
willing to move. The obstacle is visibility, not
willingness.
\begin{itemize}[leftmargin=0.7cm]
    \item \textbf{Pairwise swaps rarely work.} A direct
      exchange needs two students each to want exactly
      what the other holds. A three- or four-way
      rotation usually exists, but no participant can
      see it from inside.
    \item \textbf{No directory of intent.} Nothing
      records who wants to move or where, so willingness
      stays invisible.
    \item \textbf{Informal coordination does not scale.}
      A group chat surfaces a few offers; it cannot
      search a hostel for a closed exchange cycle.
    \item \textbf{No memory.} The same negotiation
      restarts every semester.
\end{itemize}

\textbf{Contextual relevance:} Hostel A alone has eight
floors of roughly 58 rooms. Possible exchanges grow far
faster than any student or warden can enumerate, and the
cost of missing them is borne by students for a full
academic year.

\section{Proposed Solution %
\BvrHint{\ldots Software Proposition}}
A web service that holds every current allocation and
every stated preference, and computes the exchange
cycles hidden inside them.
\begin{enumerate}[leftmargin=0.7cm]
    \item A student names one specific room. If it
      cannot be obtained, the request falls back to that
      room's floor, then its direction.
    \item The service builds a students $\times$ rooms
      cost matrix and solves for the allocation of
      maximum total satisfaction.
    \item That allocation is decomposed into independent
      exchange cycles; any cycle leaving a participant
      worse off is discarded.
    \item Each surviving cycle is issued as a proposal.
      All members approve, or it lapses.
    \item On full approval the exchange executes and new
      allocation rows are appended --- history is never
      rewritten.
\end{enumerate}
Each participant is shown their new room, the resulting
match percentage, and exactly which preferences were met
and which were not.

\section{Solution Approach %
\BvrHint{\ldots Engineering focus}}
We assemble mature engines rather than invent
algorithms. The engineering contribution is the
modelling, and the feasibility layer separating
allocations that are mathematically optimal from those
students would actually accept.

\subsection{Key modelling observation}
The pool is \emph{closed}: the only rooms available are
those the participants already hold. So any solution is
a \textbf{permutation} of students over rooms, and every
permutation decomposes uniquely into disjoint cycles.
The exchange chains are therefore not searched for
separately --- they are already inside the allocation and
are simply read out. A 2-cycle is a direct swap; a
3-cycle is a three-way rotation. Because cycles are
disjoint, each executes independently, so a participant
withdrawing from one provably cannot affect any other.

\subsection{Four modules %
\BvrHint{\ldots mature engines}}
\begin{enumerate}[leftmargin=0.7cm]
    \item \textbf{Preference scoring.} Match percentage
      is the fraction of a student's preference weight a
      room satisfies. Hard constraints rule a room out
      rather than lowering its score.
    \item \textbf{Assignment optimisation.} The
      \textbf{Hungarian algorithm}
      (\texttt{scipy.optimize.linear\_sum\_assignment})
      solves the matrix in $O(n^3)$.
    \item \textbf{Alternative generation.}
      \textbf{Murty's algorithm} yields the $K$ best
      allocations by forbidding an edge of the best and
      re-solving. This is a correctness requirement: the
      single optimal allocation is often unexecutable,
      and the alternatives are where a workable exchange
      is found.
    \item \textbf{Swap feasibility.} Cycle decomposition
      plus a Pareto filter, discarding any cycle where a
      participant ends up worse off than today.
\end{enumerate}

\subsection{Alternatives rejected}
\textbf{Gale--Shapley} assumes two-sided matching where
both sides rank each other; rooms hold no preferences
over students, so half the algorithm has no input.
\textbf{Top Trading Cycles} fits the exchange shape but
operates on a strict ranked list, and cannot represent
preferences weighted across several attributes.

\subsection{Architecture and CI/CD}
A \textbf{React} frontend over a stateless
\textbf{FastAPI} backend and managed \textbf{PostgreSQL}
(Supabase, nine tables). The allocation table is
append-only, enforced by a database trigger rather than
application convention. The domain layer imports no
framework and no database driver, so it is testable with
no infrastructure at all. GitHub Actions runs the full
suite on every push across Python 3.10--3.12 and
republishes the documentation on merge.

\section{Evaluation Criterion %
\BvrHint{\ldots measurable and attributable}}

\subsection{Primary metric}
\textbf{Cohort preference satisfaction:} the mean
weighted preference score across \emph{all} participants
--- not only those who move --- before and after a
matching round. Both figures come from the same scoring
function over the same stored preferences, so the
difference is attributable entirely to the reallocation.

On a 30-student cohort in hostel A, satisfaction rose
from \textbf{11\%} to \textbf{77\%}.

Measuring across everyone is deliberate: averaging over
movers alone would reward a two-person swap where both
reach 100\% over a 28-person rotation that lifts the
whole hostel.

\subsection{Secondary metrics}
\begin{table}[h]
\centering
\begin{tabular}{lr}
\toprule
\textbf{Metric} & \textbf{Result} \\
\midrule
Received the exact room they named & 20 of 30 \\
Received at least their requested floor & 28 of 30 \\
Received at least their requested direction & 23 of 30 \\
Participants left worse off & \textbf{0 of 30} \\
Longest chain needing simultaneous approval & 11 \\
\bottomrule
\end{tabular}
\end{table}

\textbf{Nobody worse off} is a guarantee, not an
observation: a cycle is only offered if every member is
at least as well off as before and one is strictly
better. \textbf{Chain length} matters operationally ---
an exchange needing 28 simultaneous approvals is far
more fragile than three of 11, 10 and 7. Raising $K$ cut
the longest chain from 28 to 11 \emph{at identical
satisfaction}, which is precisely what Murty's algorithm
contributes.

\subsection{Pilot validation}
Collect real preferences from 20--40 volunteers in one
hostel, run a matching round, and report before/after
satisfaction against an honest baseline: how many direct
pairwise swaps those same students could arrange
informally.

\section{Scalability %
\BvrHint{\ldots with foresight}}
\begin{enumerate}[leftmargin=0.7cm]
    \item \textbf{Theoretically:} Hungarian is $O(n^3)$
      and Murty's adds a bounded number of further
      solves. A 30-student pool matches in under
      100~ms, a 100-student pool in under a second.
      Matching also partitions naturally --- each hostel
      solved independently --- so problem size stays
      bounded as campus participation grows.
    \item \textbf{Operationally deployable:} the API is
      stateless over managed PostgreSQL, so it scales
      horizontally behind a load balancer on a standard
      platform host, with CI/CD producing repeatable
      deployments.
\end{enumerate}

\section{Engine availability heuristic}
Every component exists as a mature engine:
\textbf{SciPy} (Hungarian), \textbf{NumPy} (cost
matrix), \textbf{React} (frontend), \textbf{FastAPI}
(REST API), \textbf{Supabase} (managed PostgreSQL and
auth), \textbf{pytest} and \textbf{GitHub Actions}
(tests and CI). Murty's algorithm is about sixty lines
on top of the SciPy solver and introduces no new
mathematics --- it repeatedly forbids one assignment edge
and re-solves. Nothing here depends on unpublished or
unproven work.

\section{Project scope and deliverables %
\BvrHint{\ldots iteration-friendly}}
\subsection{Initial deliverable}
\begin{itemize}[leftmargin=0.7cm]
    \item Weighted preference model with ordered
      fallbacks: room, then floor, then direction.
    \item Matching engine: cost matrix, Hungarian,
      Murty's $K$-best, cycle decomposition, Pareto
      filter.
    \item Approval state machine: proposal, approval,
      execution, rejection, expiry, withdrawal.
    \item Web view showing the allocation table, the
      computed match, and the cycles being executed.
    \item Schema with append-only allocation history.
    \item CI running the suite on three Python versions
      on every push.
\end{itemize}

\subsection{Subsequent deliverables}
\begin{itemize}[leftmargin=0.7cm]
    \item React frontend with student login and real
      preference collection, replacing the seeded
      demonstration cohort.
    \item Persisted proposals and approvals surviving
      restart.
    \item Shared rooms, with mutual roommate requests
      merged before solving.
    \item Warden view for oversight and for committing
      an approved round.
\end{itemize}

\section{Risks and mitigations}
\begin{itemize}[leftmargin=0.7cm]
    \item \textbf{Participation:} cycles need enough
      participants. Run per hostel, where students
      already coordinate, and report honestly when no
      exchange exists.
    \item \textbf{Withdrawal mid-approval:} mitigated
      structurally --- cycles are disjoint, so a
      withdrawal cancels one and cannot touch the
      others; preferring several short cycles over one
      long one shrinks the blast radius further.
    \item \textbf{Gaming:} hard constraints only remove
      infeasible rooms and earn no score, so
      exaggerating them cannot improve a match.
    \item \textbf{Administrative acceptance:} exchanges
      remain proposals until approved, and the
      append-only history is a complete, non-rewritable
      audit trail.
    \item \textbf{Privacy:} preferences and allocations
      only; nothing beyond what matching requires.
\end{itemize}

\section{Summary}
This project proposes a deployable web service that
finds multi-party room exchanges no participant could
discover alone. It meets the course criteria through
time-to-value --- a working, measured engine in the first
iteration --- CI/CD that keeps feedback at seconds rather
than days, and a measurable evaluation criterion in
cohort preference satisfaction, which rose from 11\% to
77\% with no participant left worse off. It uses only
mature engines; the engineering contribution is
modelling the hostel as a closed permutation problem and
the feasibility layer that separates optimal allocations
from acceptable ones.

\end{document}
